% Contents of pandoc-header-pdf.tex

% Copyright Hendrik Lehmbruch 2021, wtfpl 2.0
% Copyright Axel Konrad 2021, wtfpl 2.0
% Copyright Marcel Ilg 2022, wtfpl 2.0
% see http://www.wtfpl.net/txt/copying/

% german language support
% you need to install texlive-lang-german
% comment it out, if compiling for another languge
% maybe you need another lang package (apt search latex babel)
% for LANG=C (en) you don't need an aditional package
% \usepackage[english,spanish,swedish,portuges,ngerman]{babel}
% https://www.namsu.de/Extra/papt search latex babelakete/Babel_V2017.html

\usepackage{fontenc}
\usepackage{geometry}

\usepackage{titling}
   
    % adding a logo to the title
    \pretitle{\begin{center}
    % change path to loge if necessary!
    \includegraphics[width=111.9mm,height=22.3mm]{../sys-images/logo-manual.png}\vspace{4em}\Huge\\}
    \posttitle{\end{center}\vspace{0.8em}}
    
    % author and date style
    \preauthor{\begin{center}\LARGE
    \begin{tabular}[t]{c}}
    \postauthor{\end{tabular}\par\end{center}\vspace{1em}}
    \predate{\begin{center}\large}
    \postdate{\par\end{center}}

% start sections and subsections on a new page
\AddToHook{cmd/section/before}{\clearpage}
\AddToHook{cmd/subsection/before}{\clearpage}

% header and footer
\usepackage{fancyhdr}
    % extra side for the table of contents
    \fancypagestyle{toc}
    \fancyhf{}
    \fancyhead[Re]{\thepage}
    \fancyhead[Re]{\textsl{TABLE OF CONTENTS}}
    \fancyfoot[C]{\thepage}

    % define header and footer
    \pagestyle{fancy}
    \fancyhead[C]{}
    \fancyhead[L]{\textbf{siduction Manual}}
    \fancyhead[R]{\textsl{\rightmark}}
    \fancyfoot[L,R]{}
    \fancyfoot[C]{\thepage}
    \renewcommand{\headrulewidth}{0.4pt}
    \renewcommand{\footrulewidth}{0.4pt}
    
% Paragraphs (Hurenkind und Schusterjungen verhindern)
\usepackage{indentfirst}
    \frenchspacing
\usepackage[nodisplayskipstretch]{setspace}
    \onehalfspacing
    \sloppy
    \clubpenalty10000
    \widowpenalty10000
    \displaywidowpenalty=10000

% code-block style
\usepackage{xcolor}
\usepackage{MnSymbol}% benoetigt fuer 'prebreak='
\lstset{
    basicstyle=\ttfamily\footnotesize, 
    % fontsize definitions: scriptsize,small,tiny,normalsize
    % codeline numbering
    %numbers=left,
    %numberstyle=\scriptsize,
    %stepnumber=1,
    %numbersep=8pt,
    % code highlight
    %language=bash,
    %basicstyle=\color[RGB]{40,1,196},
    %keywordstyle=\color[rgb]{0.13,0.29,0.53}\bfseries,
    %keywordstyle=\color[RGB]{42,161,152},
    %stringstyle=\color[rgb]{0.31,0.60,0.02},
    %stringstyle=\color[RGB]{28,139,210},
    %numberstyle=\color[RGB]{108,113,196},
    %identifierstyle=\color[RGB]{0,0,74},
    %commentstyle=\color[rgb]{0.56,0.35,0.01}\itshape,
    %commentstyle=\color[RGB]{0,85,0},
    %rulecolor=\color{red},
    frame=single,
    framerule=0pt,
    framexleftmargin=3pt,
    framexrightmargin=5pt,
    framextopmargin=2pt,
    framexbottommargin=2pt,
    backgroundcolor=\color[RGB]{230,230,230},
    showspaces=false,
    showstringspaces=false,
    showtabs=false,
    tabsize=2,
    captionpos=b,
    breaklines=true,
    prebreak={\raisebox{0ex}[0ex][0ex]
        {\ensuremath{\rhookswarrow}}},
    %breakatwhitespace=true,
    breakautoindent=true,
    linewidth=.95\linewidth,
    xleftmargin=14pt,
    %xrightmargin=35pt,
}

\usepackage{tcolorbox}
    \newtcolorbox{myquote}{colback=red!5!white, colframe=red!75!black}
    \renewenvironment{quote}{\begin{myquote}}{\end{myquote}}
